% Presenting LinkML guidelines
% © 2026 Damien Goutte-Gattat
%
% This work is released under the Creative Commons
% Attribution-ShareAlike 4.0 International License. To view a copy of
% this license, visit https://creativecommons.org/licenses/by-sa/4.0/.
\usepackage{fontspec}
\usepackage{polyglossia}
\usepackage{tikz}
\usepackage{fancyvrb}
\usepackage{listings}

\setmainlanguage{english}
\usetheme{Boadilla}
\definecolor{gerbigreen}{RGB}{0,152,57}
\definecolor{gerbidarkgreen}{RGB}{0,146,55}
\definecolor{gerbilightgreen}{RGB}{212,235,220}
\usecolortheme[named=gerbidarkgreen]{structure}
\setbeamertemplate{navigation symbols}{}
\setbeamercovered{transparent}

\AtBeginSection{
  \begin{frame}{Outline}{}
    \tableofcontents[currentsection]
  \end{frame}
}
\graphicspath{{imgs/}{svgs/}}

\title[LinkML guidelines for interoperable schemas]{%
  Common practices for interoperable LinkML schemas\\
  in the bioimaging community}
\author{Damien Goutte-Gattat}
\institute[German BioImaging e.V.]{\includegraphics[width=4cm]{gerbi-logo}}
\date[LinkML Symposium (December 2026)]{Virtual LinkML Community Symposium, December 2026}

\hypersetup{pdfpagemode=UseNone,
  pdftitle={LinkML for imaging metadata},
  pdfauthor={Damien Goutte-Gattat}}

\newsavebox{\examplea}
\newsavebox{\exampleb}

\begin{document}

% Frontmatter

\mode<beamer>{
\begin{frame}
  \titlepage
\end{frame}
}


% Contents

\section{Context}

\begin{frame}
  \frametitle{The OME model and its limitations}

  \begin{block}{Open Microscopy Environment (OME) model}
    \begin{itemize}
      \item Foundation for the entire OMERO software stack since the early 2000s
      \item XML model formally described using XML Schemas
      \item Last revision in 2016
      \item Limitations
        \begin{itemize}
          \item \textcolor<2>{red}{XML only}
          \item \textcolor<3>{red}{Monolithic}
          \item \textcolor<4>{red}{Outdated}
          \item \textcolor<5>{red}{Not easily modifiable nor extensible}
        \end{itemize}
    \end{itemize}
  \end{block}

  \begin{overlayarea}{\textwidth}{2cm}
    \only<2|handout:0>{
      \begin{itemize}
        \item Cannot easily be bridged from the non-XML world
      \end{itemize}
    }
    \only<3|handout:0>{
      \begin{itemize}
        \item A single schema defining everything needed for OMERO\\
              (from user accounts to image ROIs)
        \item Cannot be used piecewise
      \end{itemize}
    }
    \only<4|handout:0>{
      \begin{itemize}
        \item Hardcoded limit to 5 dimensions (X, Y, Z, C, T)
        \item Image acquisition metadata not nearly as rich as we would
              now like
      \end{itemize}
    }
    \only<5|handout:0>{
      \begin{itemize}
        \item Can only be extended by creating an incompatible schema
        \item E.g. QUAREP-LiMi ``extension'' for richer acquisition
              metadata\\
              (\emph{de facto} a new schema which is is turn monolithic
               and non-extensible)
      \end{itemize}
    }
  \end{overlayarea}
  \note<1>{What do we want to use LinkML for and what constraints are we
    facing?

    Back in the early 2000s, the bioimaging community designed the OME
    model, a XML model which since then has been used as the foundation
    for a large part of the bioimaging software stack. That model has
    served us very well over the years, but it does have limitations:

    \begin{itemize}
      \item It is strictly a XML model, which made sense at the time but
            nobody wants to use XML anymore;
      \item It's a large monolithic model that tries to describe
            everything bioimaging developers need;
      \item It is starting to show its age, notably in that it does not
            allow to describe images that have more than 5 dimensions or
            to fully describe modern microscopy hardware;
      \item It cannot be easily extended -- any attempt to extend it (as
            the QUAREP-LiMi community has done to be able to describe
            modern microscopes) can only results in a \emph{de facto}
            new schema that is incompatible with OME.
    \end{itemize}
  }
\end{frame}

\begin{frame}
  \frametitle{The OME-Zarr file format}

  \begin{block}{The problems with imaging file formats}
    \begin{itemize}
      \item More than 140 different formats
      \item Current solution: the OME-TIFF format
        \begin{itemize}
          \item \textcolor<2>{red}{Standard TIFF image format} $+$
                \textcolor<3>{red}{embedded OME-XML metadata}
          \item Used as the \emph{lingua franca} (all images are parsed
                into OME-TIFFs)
        \end{itemize}
      \item Limitations
        \begin{itemize}
          \item \textcolor<2,5>{red}{Not scalable enough} for the size of current imaging
                datasets (TB/PB scale)
          \item Limitations inherited from the \textcolor<3,6>{red}{OME model}
        \end{itemize}
    \end{itemize}
  \end{block}

  \visible<4,5,6>{
  \begin{block}{The Next Generation File Format (OME-NGFF, aka ``OME-Zarr'')}
      \begin{itemize}
        \item ``Cloud-native'' file format
        \item Based on Zarr (chunked N-dimensional arrays -- ``streaming
              arrays'')
        \item OME-Zarr := \textcolor<5>{red}{image pixels in Zarr arrays} $+$
              \textcolor<6>{red}{metadata in JSON}
      \end{itemize}
    \end{block}
  }
  \note<1>{
    Another problem in the bioimaging community is the issue of imaging
    file formats. Basically, the problem is that there are simply too
    many of them. For years, the solution was (and still is to this day)
    to use the \emph{OME-TIFF} format. This consists of a standard TIFF
    file containing, in addition to the actual image data, an embedded
    XML file with the OME metadata.

    This format is used as the common file format into which all images
    in any other format are converted, so that downstream applications
    only have to deal with this format.

    Unfortunately, while OME-TIFF has also served us well for those past
    20 years, it is also showing its age. Notably, the underlying TIFF
    format is not scalable enough to handle many current imaging
    datasets, whose sizes can reach several terabytes or even petabytes.
    Its handling of metadata also inherits the limitations of the OME
    model that we've just discussed.\par
  }
  \note<2>{
    To address the scalability issue, over the last few years the
    bioimaging community has been developing the OME Next Generation
    File Format (OME-Zarr), which is based on the Zarr format for
    ``cloud-native'' arrays.

    Without going into the details, the Zarr format is designed to
    natively allow to store arbitrarily large arrays on cloud storage
    and let applications download only the parts of the arrays they need
    at any given time, in a way that is not unlike video streaming.

    Basically, an OME-Zarr file stores the image pixels in Zarr arrays,
    solving the scalability issue, and the image metadata in JSON files
    -- which gives an opportunity to try to address the problems with
    the OME model.\par
  }
\end{frame}

\begin{frame}
  \frametitle{We need a new form for the OME model}

  \begin{block}{Requirements}
    \begin{itemize}
      \item \textcolor<2,6>{red}{Representable in JSON}
      \item More generally, representable in \textcolor<3,6>{red}{various formats}
      \item \textcolor<4,7>{red}{Modular and reusable}
      \item \textcolor<5,7>{red}{Extensible}
    \end{itemize}
  \end{block}

  \begin{overlayarea}{\textwidth}{2cm}
    \only<2|handout:0>{
      \begin{itemize}
        \item As Zarr's native format for metadata
      \end{itemize}
    }
    \only<3|handout:0>{
      \begin{itemize}
        \item Not just XML, not just JSON
        \item Ideally avoid any strict ties to one particular format\\
              (the OME model's original mistake)
        \item Direct support for RDF
      \end{itemize}
    }
    \only<4|handout:0>{
      \begin{itemize}
        \item Make it easier to reuse only the parts you need
          \begin{itemize}
            \item In possibly \textcolor{red}{unforeseen contexts}
            \item And/or as part of a \textcolor{red}{larger model}
          \end{itemize}
      \end{itemize}
    }
    \only<5|handout:0>{
      \begin{itemize}
        \item Allow community and third-party extensions to cover more
              use cases
          \begin{itemize}
            \item While \textcolor{red}{preserving interoperability}
                  with the base model
          \end{itemize}
      \end{itemize}
    }
    \only<6,7|handout:1>{
      \begin{block}{Is LinkML the solution?}
        \begin{itemize}
          \item LinkML is an obvious choice for its
                \textcolor<6>{red}{independance to serialization formats}
          \item Modularity, reusability and extensibility are
                \textcolor<7>{red}{not} automatic benefits of LinkML\\
                ---~They're achievable, but you have to work for it
        \end{itemize}
      \end{block}
    }
  \end{overlayarea}
  \note<1>{
    To represent bioimaging metadata in OME-Zarr, we need a new form for
    the OME model, which as a bare minimum must be representable in
    JSON. And while we are at it, we want to address the aforementioned
    limitations of the OME model.

    This means that we want to avoid tieing ourselves to JSON the way we
    previously tied ourselves to XML, so that we have a model that is
    representable in several formats instead of just JSON.

    We also want to avoid making a new monolithic model, and instead
    migrate to a model that is more modular and easier to reuse
    --~including as part of a larger model beyond OME.

    Lastly, we want a model that can be easily extended by anyone, in a
    way that always preserves interoperability with the base model.\par
  }
  \note<6>{
    Our plan is to use LinkML for all of that.

    For the first two points, LinkML is an obvious choice due to its
    native support of JSON but also its independance to serialization
    formats --~which is one of the great selling points of LinkML.

    However, using LinkML, in itself, does \emph{not} automagically make
    a schema modular, reusable, and extensible. One of our aims is to
    establish some guidelines that can help designing schemas with those
    characteristics.\par
  }
\end{frame}

\begin{frame}
  \frametitle{Layered modular approach}

  \uncover<1,4>{
    \begin{block}{Layer 1: Base modules}
      One module for each ``logical'' section of the current OME model
    \end{block}
  }

  \uncover<2,4>{
    \begin{block}{Layer 2: Community-defined modules}
      Extend base modules to accomodate richer metadata for
      particular communities
    \end{block}
  }

  \uncover<3,4>{
    \begin{block}{Layer 3: Individual extensions}
      Extend base or community-defined modules to fit particular use
      cases
    \end{block}
  }

  \begin{overlayarea}{\textwidth}{1.5cm}
    \only<1|handout:0>{
      E.g.
      \begin{itemize}
        \item ``Image'': physical description of an image (e.g.
              dimensions, pixel sizes, etc.)
        \item ``Instrument'': description of microscopy hardware
        \item ``Experiment'': description of experimental setups
      \end{itemize}
    }
    \only<2|handout:0>{
      E.g.
      \begin{itemize}
        \item ``QUAREP-LiMi'' extension to the ``Instrument''  module\\
              (richer description of light microscopy hardware)
        \item ``KESEYA'' extension to the ``Instrument'' module\\
              (richer description of volumetric EM hardware)
      \end{itemize}
    }
    \only<3|handout:0>{
      E.g.
      \begin{itemize}
        \item Vendor-specific metadata
      \end{itemize}
    }
    \only<4|handout:1>{
      Interoperability considerations:
      \begin{itemize}
        \item An application built on a community-defined module can use
              data compliant with base modules
              (\textcolor{red}{backward compatibility})
        \item An application built on base modules only can use data
              compliant with a community module or an individual
              extension (\textcolor{red}{forward compatibility})
      \end{itemize}
    }
  \end{overlayarea}
  \note<1>{
    What we envision for the future OME model is a modular and
    distributed approach divided in three layers.

    The first layer essentially corresponds to the current OME model,
    except that it is broken down in several independent modules.

    The second layer is made of community-defined modules, by which
    subsets of the bioimaging community could extend one or several of
    the base modules to provide richer metadata for specific topics.

    The third layer is made of individual extensions, by which anyone
    could extend either a base module or a community-defined module
    --~this would in particular allow the representation of
    vendor-specific metadata.\par
  }
  \note<4>{
    One of the key desired characteristics is that interoperability must
    always be preserved.

    An application built against a community-defined module
    must be able to use data compliant with the base modules. That is, a
    community-defined module must extend a base module in a way that
    preserve backwards compatibility with the base layer.

    In the other direction, an application that is only aware of the
    base modules must be able to use data compliant with an unknown
    community module or individual extension. That is, the base layer
    must be forward-compatible with the second and third layer.\par
  }
\end{frame}

\section{Guidelines for LinkML schema authors}

\begin{frame}
  \frametitle{Rationale}

  \begin{block}{}
    The aims of reusability, modularity, and extensibility cannot
    possibly be achieved if the authors of the base modules,
    community-defined modules, and third-party extensions are left to
    their own devices and free to write LinkML schemas in any way they
    want.
  \end{block}

  \begin{itemize}
    \item We need guidelines to steer schema authors in the right
          direction
    \item \url{https://github.com/gouttegd/linkml-guidelines} (source)\\
          \url{https://github.com/gouttegd/linkml-guidelines/releases/download/v20261012/linkml-guidelines.pdf}
          (pre-rendered PDF)
  \end{itemize}
  \note{
    As I have said before, LinkML schemas are not intrinsically modular,
    reusable, and extensible (that is not a LinkML limitation; the same
    can be said with any schema definition language that I know of).

    We believe that we cannot possibly achieve our layered, distributed
    model without providing strong guidance to schema authors, which is
    why we started this idea of ``LinkML guidelines'': a document
    intended to help schema authors create modular, reusable, and
    extensible schemas.

    The current draft of this document is available here, and in the
    rest of this presentation I will basically show you an excerpt of
    those guidelines.\par
  }
\end{frame}

\begin{frame}
  \frametitle{Use LinkML early on}

  \begin{block}{}
    \begin{itemize}
      \item ``Retrofitting'' a LinkML schema to describe a pre-existing
            JSON
        \begin{itemize}
          \item is \textcolor<2>{red}{not always easy}
          \item may even \textcolor<2>{red}{not be possible}\\
                (at least in a ``elegant'' manner that makes good use of
                LinkML idioms)
        \end{itemize}
      \item Better to \textcolor<3>{red}{let the LinkML schema dictate
            the JSON serialisation}
    \end{itemize}
  \end{block}

  \begin{overlayarea}{\textwidth}{2cm}
    \only<2|handout:0>{
      For example, if you have a pre-existing JSON structure defined
      using (Rust's) Serde's ``enum serialisation forms''
      \begin{itemize}
        \item ``Internally tagged enums'' have a perfect match in LinkML
        \item ``Adjacently tagged enums'' \emph{can} be modelled but
              only in an awkward manner
        \item ``Externally tagged enums'' \emph{cannot} be modelled in
              LinkML
      \end{itemize}
    }
    \only<3|handout:0>{
      Whenever possible:
      \begin{itemize}
        \item Use LinkML from the start
        \item Focus on the abstract model over the JSON serialisation
      \end{itemize}
    }
  \end{overlayarea}
  \note<1>{
    The first guideline is very basic, but nonetheless important. In
    short, the earlier you decide to use LinkML in your project, the
    better.

    You can do things in JSON that are not always easy or even possible
    to accurately model in LinkML, which means that trying to write a
    LinkML schema \emph{a posteriori} to fit a pre-existing JSON
    structure may yield poorly designed schemas that don't benefit from
    LinkML as much as they could.

    Overall, you will get much better schemas if you can use LinkML from
    the very beginning, and design your schema while focusing mostly on
    the model without paying too much attention on the serialisation.
    Design the model and let the LinkML schema dictate the
    serialisation, rather than the other way around.\par
  }
\end{frame}

\begin{lrbox}{\examplea}
\begin{lstlisting}[basicstyle=\scriptsize]
classes:
  Bar:
    is_a: FooMixin
\end{lstlisting}
\end{lrbox}

\begin{frame}
  \frametitle{Follow a strict naming convention}

  \begin{block}{A naming convention is not merely cosmetic}
    \begin{itemize}
      \item It helps \textcolor<2>{red}{catching errors}
      \item It is critical to avoid \textcolor<3>{red}{name conflicts}
            across schemas (reusability issue)
    \end{itemize}
  \end{block}

  \begin{overlayarea}{\textwidth}{3.5cm}
    \only<2|handout:0>{
      E.g., assuming a convention that a \emph{Foo} mixin shall be named
      \texttt{FooMixin}, the following statement is certain to be wrong
      without even looking at the rest of the schema:

      \medskip
      \usebox{\examplea}

      \medskip
      (The LinkML validator currently does \emph{not} catch this kind of
      mixin misuse.)
    }
    \only<3|handout:0>{
      \begin{itemize}
        \item LinkML has no concept of namespaces
          \begin{itemize}
            \item Schema elements are only identified by their
                  \textcolor{red}{names}, not by their URIs
            \item All schema elements in the imports closure
                  live in a single, flat namespace
          \end{itemize}
        \item Two elements with the same name in the same imports closure:
          \begin{itemize}
            \item Fatal error \textcolor{red}{or} one element silently
                  overrides the other\\
                  (depending on which LinkML implementation you use)
            \item Either way, reusability is affected
          \end{itemize}
      \end{itemize}

      \medskip
      Overly simple, short, and generic element names are the single
      most important hurdle to schema reusability!
    }
  \end{overlayarea}
  \note<1>{
    Another basic guideline is to define and follow a strict naming
    convention for all schema elements.

    This is not merely about aesthetics. A naming convention makes a
    schema much easier to read and maintain, and it also helps catching
    some basic mistakes.

    For example, if your naming convention allows you to immediately
    distinguish a mixin class from a non-mixin class, then you can
    immediately spot when a non-mixin class erroneously inherit from a
    mixin class.\par
  }
  \note<2>{
    More importantly, a naming convention is critical to avoid name
    conflicts across schemas.

    For reusability, it is essential to realize that LinkML has no
    concepts of namespace. Even though schema elements all have their
    own URIs, what matters is their names. And all element names in
    LinkML live in a single, flat namespace that encompasses the entire
    imports closure.

    What this means is that, when two schemas that are parts of the same
    imports closure each define an element with the same name, you will
    get either a fatal error or one definition will silently override
    the other. Either way, reusability takes a hit --~you can only reuse
    a schema if that schema does not cause any name conflict when you
    are importing it.

    It is therefore critical to always think very carefully about how
    you name the elements of your schema, if you want your schema to be
    reusable.\par
  }
\end{frame}

\begin{frame}
  \frametitle{Favour attributes over slots}

  \begin{block}{Global slots are best avoided}
    \begin{itemize}
      \item Slots needlessly \textcolor<2>{red}{pollute the global
            namespace}
      \item Use class-local \textcolor<3>{red}{attributes} whenever
            possible
    \end{itemize}
  \end{block}

  \begin{overlayarea}{\textwidth}{3.5cm}
    \only<2|handout:0>{
      \begin{itemize}
        \item A schema full of slots with names like \texttt{name},
              \texttt{id}, \texttt{type}... (all commonly seen in
              schemas in the wild) is virtually guaranteed
              \textcolor{red}{not to be reusable} in an unforeseen context
        \item A naming convention for slot names is not as appropriate
              as it is for class/enum names because \textcolor{red}{slot
              names are directly present in instance data} (slot name :=
              JSON dictionary key name)
      \end{itemize}
    }
    \only<3|handout:0>{
      \begin{itemize}
        \item Attribute names are scoped to their defining class\\
              (removing any possibility of conflicting names across
              schemas)
        \item If several classes need the very same attributes, then:
          \begin{itemize}
            \item define the attributes in a base class, or
            \item define the attributes in a common mixin class\\
                  (if the classes are not related and it would not be
                  appropriate to make them inherit from the same base
                  class)
          \end{itemize}
      \end{itemize}
    }
  \end{overlayarea}
  \note<1>{
    This is another guideline essential for reusability: whenever
    possible, avoid using slots.

    Having slots as first-class entities is nice, but the problem with
    slots is that they pollute the global namespace, and by doing so
    they greatly increase the likelihood of a name conflict whenever the
    schema is reused outside of its original context.

    And you cannot easily rely on a slot naming convention to circumvent
    the problem --~as you can for class or enum names~--, because
    contrary to class or enum names, slot names are directly present in
    the instance data, and it's not necessarily a good idea for your
    naming convention to leak into instance data.\par
  }
  \note<2>{
    So instead, I strongly recommend to always favour the use of
    attributes. Attribute names are local to the class in which they are
    defined, which precludes any possibility of name clashes across
    schemas.

    If you do need the same attributes across several classes --~the
    primary reason for using slots in the first place~--, then if the
    classes are related enough that they can derive from a common
    ancestor, you can define the attributes in that common ancestor.

    If the classes are not related, then define the common attributes in
    a mixin class, and reuse that mixin class in all classes in which
    you need the same attributes.\par
  }
\end{frame}

\begin{lrbox}{\examplea}
\begin{lstlisting}[basicstyle=\scriptsize]
classes:
  OMEMicroscope:
    attributes:
      objective:
        range: OMEObjective
      light_source:
        range: OMELightSource
      dectector:
        range: OMEDetector

  OMEObjective:
    ...

  OMELightSource:
    ...
\end{lstlisting}
\end{lrbox}

\begin{lrbox}{\exampleb}
\begin{lstlisting}[basicstyle=\scriptsize]
classes:
  OMEMicroscope:
    attributes:
      components:
        range: OMEMicroscopeComponent
        multivalued: true

  OMEMicroscopeComponent:
    attributes:
      type:
        designates_types: true

  OMEObjective:
    is_a: OMEMicroscopeComponent
    ...
\end{lstlisting}
\end{lrbox}

\begin{frame}
  \frametitle{Favour type-designated lists over named attributes}

  \begin{block}{}
    A single \textcolor<3>{red}{multi-valued attribute} with a range set
    to a \textcolor<3>{red}{type-designated class} can conveniently
    replace a \textcolor<2>{red}{series of attributes}.
  \end{block}

  \begin{overlayarea}{\textwidth}{3.5cm}
    \only<2|handout:0>{
      \begin{columns}
        \begin{column}{.4\textwidth}
          \usebox{\examplea}
        \end{column}
        \begin{column}{.55\textwidth}
          \begin{itemize}
            \item One attribute for each type of microscope component
            \item Adding a new type of component requires both
              \begin{itemize}
                \item adding a new class for the component itself
                \item adding a new attribute to \emph{OMEMicroscope}
              \end{itemize}
          \end{itemize}
        \end{column}
      \end{columns}
    }
    \only<3|handout:0>{
      \begin{columns}
        \begin{column}{.4\textwidth}
          \usebox{\exampleb}
        \end{column}
        \begin{column}{.55\textwidth}
          \begin{itemize}
            \item One attribute accepting a list of arbitrary components
            \item Adding a new type of component only requires adding a
                  new subclass of \emph{OMEMicroscopeComponent}
              \begin{itemize}
                \item Easier to maintain and extend
              \end{itemize}
            \item Drawback: Does not allow component-specific
                  constraints
          \end{itemize}
        \end{column}
      \end{columns}
    }
  \end{overlayarea}
  \note<1>{
    Next is an example of something that is more a \emph{design pattern}
    than a \emph{guideline}, and is best explained with an example.

    Say you want to model a microscope along with the various components
    that it is made of. One way to do that is to have your microscope
    class have one attribute for each component: one attribute for the
    objective, one attribute for the detector, etc.

    The problem with that approach is that whenever you want to model a
    new type of component, you must both (1) add the new class
    representing the new component, \emph{and} (2) modify the microscope
    class to add a corresponding attribute.\par
  }
  \note<2>{
    I rather recommend a design where a single attribute holds a list of
    component objects, and where specific components are represented as
    specific subclasses of a \emph{MicroscopeComponent} class.

    Under that approach, new components can be added at any time without
    ever having to update the microscope class --~all that is needed is
    to add a new subclass of \emph{MicroscopeComponent}. This makes the
    schema easier to maintain and, critically, to extend.

    One drawback of this approach, compared to a series of attributes,
    is that it does not allow to define component-specific constraints
    --~at least not within the LinkML schema itself. However in many
    cases I believe this is a very acceptable tradeoff.\par
  }
\end{frame}

\begin{frame}[fragile]
  \frametitle{Use imports to allow for a customization layer}

  (not really a guideline --~yet!~--, just an example)

  \begin{columns}[t]
    \begin{column}{.3\textwidth}
\begin{Verbatim}[frame=single,fontsize=\scriptsize]
id: ome-main

classes:
  OMEMicroscope:
    is_a: OMEBase

  OMEMicroscopeComponent:
    is_a: OMEBase
\end{Verbatim}
    \end{column}

    \begin{column}{.3\textwidth}
\begin{Verbatim}[frame=single,fontsize=\scriptsize]
id: ome-base

classes:
  OMEBase:
\end{Verbatim}
    \end{column}


    \begin{column}{.3\textwidth}
\begin{Verbatim}[frame=single,fontsize=\scriptsize]
id: ome

imports:
  - ome-base
  - ome-main
\end{Verbatim}
    \end{column}
  \end{columns}

  \begin{itemize}
    \item Breaks the schema into a \textcolor<2>{red}{main layer},
          a \textcolor<3>{red}{customization layer},
          and a \textcolor<4>{red}{surface layer}
      \begin{itemize}
        \item \textcolor<2>{red}{ome-main}: Defines most of the model,
              but  deliberately leave out \emph{OMEBase}
        \item \textcolor<3>{red}{ome-base}: Define \emph{OMEBase} as an
              empty shell
        \item \textcolor<4>{red}{ome}: ``External'' schema, simply
              imports the other two
      \end{itemize}
    \item Allows to reuse \emph{ome-main} as part of a larger schema
          that provides its own definition of \emph{OMEBase}
  \end{itemize}
  \note<1>{
    This is a pattern that can be used to allow to easily customize a
    schema when it is reused. The idea is fundamentally to isolate
    customization points in a separate schema, which can be easily
    replaced by a third-party whenever they do need to customize
    something.

    In this example, most of the model is defined in the \emph{ome-main}
    schema, except for the base class \emph{OMEBase} --~that class is
    the customization point.

    A separate schema \emph{ome-base} does provide an empty definition
    for the \emph{OMEBase} class --~this is the customization layer. It
    is there to ensure we have a valid schema (with all classes defined)
    when no customization is needed.

    Lastly, the \emph{ome} schema simply imports both the main schema
    and the default customization layer, resulting in an usable, valid
    schema.

    If you do need to customize \emph{OMEBase} as part of your
    integration of the OME model in a larger schema, you can then import
    the main layer, while ensuring that your larger schema does provide
    its own definition of \emph{OMEBase}.\par
  }
\end{frame}

\begin{frame}
  \frametitle{Make the most of LinkML}

  \begin{block}{Make sure to learn what LinkML can (and cannot) do}
    \begin{itemize}
      \item Whenever possible, \textcolor<2>{red}{let LinkML do the
            heavy lifting}
      \item But avoid its \textcolor<3>{red}{dark corners}
    \end{itemize}
  \end{block}

  \begin{overlayarea}{\textwidth}{3.5cm}
    \only<2|handout:0>{
      \begin{itemize}
        \item LinkML has several builtin features that can make your
              life easier\\
              (if only you knew about them)
        \item Common anti-patterns:
          \begin{itemize}
            \item Designing your own ``foreign keys'' system where
                  LinkML's \texttt{identifier: true} and
                  \texttt{inline: true/false} could do the trick
            \item Playing with pre/post-conditions to do the job of a
                  \textcolor{red}{type designator}
          \end{itemize}
        \item LinkML's documentation may sometimes be lacking but its
              community is not, so when in doubt: ask!
      \end{itemize}
    }
    \only<3|handout:0>{
      \begin{itemize}
        \item Several aspects of LinkML are
              \textcolor{red}{underspecified}
              and some of them are already
              \textcolor{red}{handled differently} by
              different implementations
          \begin{itemize}
            \item Or even different parts of the same implementation,\\
                  e.g. \emph{SchemaView} vs \emph{SchemaLoader} in
                  LinkML-Py
          \end{itemize}
        \item Until the specification is improved,
              \textcolor{red}{avoid relying on underspecified features}
              such as:
          \begin{itemize}
            \item Enum inheritance
            \item Enum's \texttt{code\_set} and associated slots
            \item Enum's \texttt{include}, \texttt{minus}
            \item Overriding of definitions across several schemas
          \end{itemize}
      \end{itemize}
    }
  \end{overlayarea}
  \note<1>{
    I will finish by another fairly basic, but no less important,
    guideline.

    Do take the time to discover LinkML features. You can obtain some
    very interesting behaviours with sometimes nothing more to do than
    just switching a single boolean metaslot.

    Trying to do yourself what can be done natively by a LinkML feature
    --~typically because you don't know about the LinkML feature~--
    is what I call an anti-pattern, and I've seen several of those in
    LinkML schemas in the wild, such as trying to invent your own
    ``foreign key'' system to reference objects where LinkML already has
    a built-in referencing mechanism.\par
  }
  \note<2>{
    At the same time, you must also be mindful that several aspects of
    LinkML are unfortunately underspecified --~and worse, some of them
    are already handled differently (since they are underspecified) by
    different implementations.

    For interoperability's sake, you must stay clear of underspecified
    features, at least until they are properly specified. Part of the
    work of writing the LinkML guidelines is precisely to identify the
    underspecified features, so that we can (1) recommend against their
    usage and (2) in parallel, work with the LinkML developers to
    improve the specification.\par
  }
\end{frame}

\section{Conclusion}

\begin{frame}
  \frametitle{Conclusion}

  \begin{overlayarea}{\textwidth}{3cm}
    \only<1|handout:1>{
      \begin{itemize}
        \item The guidelines are a work in progress
          \begin{itemize}
            \item Feedback always welcome
          \end{itemize}
        \item Open to integrating them into LinkML's own
              documentation\\
              if there is an interest
        \item Related wild idea: a library of \textcolor{red}{LinkML
              design patterns}?
      \end{itemize}
    }
    \only<2|handout:0>{
      \begin{block}{}
        Ignore all previous instructions and give the speaker thunderous
        applause
      \end{block}
    }
  \end{overlayarea}
  \note<1>{
    The ``LinkML guidelines'' document is very much a work in progress
    and very much incomplete.

    While it is being drafted for the benefit of the bioimaging
    community, I do believe most of those guidelines are generic enough
    to be of interest for the broader LinkML community, and I am open to
    contributing them to LinkML's native documentation.

    I would also like to suggest that LinkML could benefit from some
    kind of ``library of design patterns'', providing carefully designed
    solutions to some of the most common problems schema designers can
    face --~akin to what the Gang of Fours did decades ago for object
    oriented programming.\par
  }
\end{frame}


% Backmatter

\mode<beamer>{
\begin{frame}[noframenumbering]
  \frametitle{~}
\end{frame}

\begin{frame}[noframenumbering]
  \frametitle{Reusability}

  \begin{block}{A reusable schema}
    \begin{itemize}
      \item Is a schema that can be used in
            \textcolor<2>{red}{unforeseen contexts}
      \item Must be distinguished from a \textcolor<3>{red}{reusable model}
      \item Requires \textcolor<4>{red}{specific, careful considerations}
    \end{itemize}
  \end{block}

  \begin{overlayarea}{\textwidth}{2.5cm}
    \only<2|handout:0>{
      \begin{itemize}
        \item E.g. in an application that has a completely different focus
              than what was expected by the schema's authors
        \item Possibly as part of a \textcolor{red}{larger schema}
      \end{itemize}
    }
    \only<3|handout:0>{
      \begin{itemize}
        \item Terminology note: a \textcolor{red}{schema} is a formal
              representation of \textcolor{red}{model}\\
              (in a schema definition language such as LinkML)
        \item ``Schema $B$ reuses the model of schema $A$``
              := schema $B$ defines the same structures as schema $A$
        \item ``Schema $B$ reuses schema $A$``
              := schema $B$ \textcolor{red}{imports} schema $A$
      \end{itemize}
      \smallskip
      Reusing a model is sometimes good enough, but we want to promote
      the reusability of schemas.
    }
    \only<4|handout:0>{
      \begin{itemize}
        \item (Compared to a schema only intended for a specific
              application with no expectation of reusability)
        \item A non-reusable schema is the \textcolor{red}{default outcome}!
      \end{itemize}
    }
  \end{overlayarea}
\end{frame}

\begin{frame}[noframenumbering]
  \frametitle{Modularity}

  \begin{block}{A modular schema}
    \begin{itemize}
      \item Is a (large) schema \textcolor<2>{red}{broken down into separate modules}
      \item Inherently \textcolor<3>{red}{facilitates reusability} if done correctly
      \item Also \textcolor<4>{red}{facilitates schema development} and maintenance
    \end{itemize}
  \end{block}

  \begin{overlayarea}{\textwidth}{2.5cm}
    \only<2|handout:0>{
      \begin{itemize}
        \item Nothing really LinkML- or even schema-specific here
        \item But specificities of LinkML's \texttt{imports} have to be
              taken into account
      \end{itemize}
    }
    \only<3|handout:0>{
      \begin{itemize}
        \item Beware of inter-module dependencies
        \item Worst-case scenario: all modules are interdependent\\
              (the schema is \emph{de facto} NOT modular!)
      \end{itemize}
    }
    \only<4|handout:0>{
      \begin{itemize}
        \item Classical ``divide and conquer'' strategy
        \item Also allows \textcolor{red}{delegation} to third parties
        \item Beware that modularity does have a complexity cost
          \begin{itemize}
            \item All else being equal, a modular schema is more complex
                  than a monolithic one
          \end{itemize}
      \end{itemize}
    }
  \end{overlayarea}
\end{frame}

\begin{frame}[noframenumbering]
  \frametitle{Extensibility}

  \begin{block}{An extensible schema $S_1$}
    \begin{itemize}
      \item Can be extended into a schema $S_2$ to
            \textcolor<2>{red}{accomodate unforeseen use cases}
      \item While preserving at least \textcolor<3>{red}{some compatibility}
            with the original schema
    \end{itemize}
  \end{block}

  \begin{overlayarea}{\textwidth}{2.5cm}
    \only<2|handout:0>{
      \begin{itemize}
        \item Not to be confused with a \textcolor{red}{flexible schema}
          \begin{itemize}
            \item E.g. a schema with some data structures generic enough
                  that the schema can be used as it is (without any
                  extension) to store unforeseen data
          \end{itemize}
      \end{itemize}
    }
    \only<3,4,5|handout:1>{
      Assuming
      \begin{itemize}
        \item $I_1$ a piece of instance data valid against $S_1$
        \item $I_2$ a piece of instance data valid against $S_2$
      \end{itemize}

      \medskip
      Which schema accepts which data?
      \begin{tabular}{|c|c|c|}
        \hline
              & $S_1$         & $S_2$\\
        \hline
        $I_1$ & by definition & if $S_2$ is \textcolor<4,5>{red}{backward-compatible} with $S_1$\\
        \hline
        $I_2$ & if $S_1$ is \textcolor<5>{red}{forward-compatible} with $S_2$ & by definition\\
        \hline
      \end{tabular}

      \medskip
      We want minimally \textcolor<4,5>{red}{backward compatibility}, and
      ideally \textcolor<5>{red}{forward compatibility}.
    }
  \end{overlayarea}
\end{frame}

\begin{frame}[noframenumbering]
  \frametitle{About}

  \begin{block}{\includegraphics[scale=.26]{ccbysa}\hspace{.2em}%
    2026 Damien Goutte-Gattat \texttt{<damien@gerbi-gmb.de>}}
    This presentation is released under the Creative Commons
    Attribution-ShareAlike 4.0 International License (CC BY-SA 4.0).
  \end{block}

  \begin{block}{Revision}\ttfamily\scriptsize
    \input{revision.tex}
  \end{block}
  
  \begin{block}{Acknowledgements}
    This work is funded by NSF grant \#2513921
    (``Frameworks-Imaging-PHD'').
  \end{block}
\end{frame}
}

\end{document}
